% Corte mecánico íntegro: parte_i.tex:420--528 + parte_ii.tex:1--321
% (la línea 322 del propietario es un separador vacío).
% Residencia material del ensamblaje sucesor de 102 capítulos.
% Residencia causal: capítulo 8; no sustituye al propietario Γ_9.
\chapter{TPK: selección, transporte, actualización y retorno}

\section{La transición tipada}

El TPK no es un contenedor nominal. La función de fase y su vector tabulado
son, respectivamente,
\[
 M_{\ph}:\mathcal D_9\longrightarrow\{+1,-1,0\},
 \qquad
 m_{\ph}:=\bigl(M_{\ph}(1),\ldots,M_{\ph}(9)\bigr)^{\mathsf T}.
\]
No son el mismo tipo de objeto y ninguno de ellos actúa por composición
directa sobre el estado. Para \(x\in X_{\TPK}^{\enr}\), definimos
\[
 \operatorname{Tra}_t(y,d):=\operatorname{Lift}(T_d)(y),
 \qquad
 \operatorname{Upd}_t:=\operatorname{Rec}_t\circ\operatorname{Mem}_t.
\]
La transición elemental se escribe entonces
\begin{equation}
 U_t(x)=\operatorname{Upd}_t\!\left(
 \operatorname{Tra}_t\!\left(
 \operatorname{Sel}_t\!\left(x,M_{\ph}(g(t))\right),d(t)
 \right)\right),
 \qquad
 U_t:X_{\TPK}^{\enr}\longrightarrow X_{\TPK}^{\enr}.
 \label{espiral:eq:transicion-tpk-tipada}
\end{equation}
El selector \(\operatorname{Sel}_t\) lee el valor escalar del régimen
trítico y elige hoja aditiva, multiplicativa o umbral;
\(\operatorname{Tra}_t(\,\cdot\,,d(t))=\operatorname{Lift}(T_{d(t)})\)
aplica la dirección cardinal;
\(\operatorname{Mem}_t\) actualiza residuo, cociente, acarreo, firma,
cilindro, frontera y memoria; y \(\operatorname{Rec}_t\) efectúa el retorno
tipado. La escritura abreviada
\(\mathrm{Upd}_t\circ\mathrm{Tra}_t\circ\mathrm{Sel}_t\) designa esta misma
composición y no un segundo motor.

La fibra enriquecida mínima puede representarse como
\[
 \mathcal F_q=O_q\times H_q\times C_q\times S_q\times B_q\times\Lambda_q,
\]
con orientación/hoja, elevación, acarreo, firma, frontera y registro. El
detalle de cada factor depende de la residencia, pero la regla de integridad
es estable: una coordenada no puede eliminarse si modifica una continuación.

\section{Holonomía nonádica}

Para un bloque compatible de nueve pasos,
\begin{equation}
 \Gamma_{9,K}=U_{K+8}\circ\cdots\circ U_K.
 \label{espiral:eq:gamma9-def}
\end{equation}
La fase satisface
\[
 g(K+9)=g(K),
\]
pero el estado enriquecido no retorna en general:
\[
 \Gamma_{9,K}(x_K)\neq x_K.
\]
En particular,
\[
 q_{\mem}(\Gamma_9x)=q_{\mem}(x)+1,
 \qquad B_{K+9}\neq B_K\quad\text{en general}.
\]

\begin{theorem}[Retorno de fase con avance de memoria]
La holonomía \(\Gamma_9\) cierra la coordenada de fase y transporta el
estado a una nueva profundidad. Por tanto, el retorno de fase no es retorno
de estado.
\end{theorem}
\begin{proof}
El cierre de fase es la periodicidad de \(g\). La actualización del ledger
incluye el incremento de profundidad y el transporte de frontera. Si el
estado retornara, todas las coordenadas deberían coincidir, contradiciendo
\(q_{\mem}\mapsto q_{\mem}+1\). Por tanto sólo retorna la fase.
\end{proof}

\section{Composición del acarreo}

Sea \(\mathcal K_{\Carr}\) el módulo de acarreo y sean
\(H,Y:\mathscr P_{\TPK}^{\enr}\to\operatorname{End}(\mathcal K_{\Carr})\)
las acciones derecha e izquierda inducidas, respectivamente, por el prefijo
conservado y por la hoja activa de una ruta. Consideramos
\(\mathcal K_{\Carr}\) como bimódulo y escribimos
\(kH(\gamma)\) y \(Y(\gamma)k\) para esas acciones. El mapa
\(\Carr:\mathscr P_{\TPK}^{\enr}\to\mathcal K_{\Carr}\) asigna a cada ruta
su acarreo total.

Para rutas componibles \(\gamma_1,\gamma_2\), el acarreo cumple
\begin{equation}
 \Carr(\gamma_2\gamma_1)
 =\Carr(\gamma_2)H(\gamma_1)
 +Y(\gamma_2)\Carr(\gamma_1).
 \label{espiral:eq:cociclo-carry}
\end{equation}
Esta ley muestra qué significa conservar memoria: no mantener constante una
cantidad, sino conservar su regla exacta de actualización. La memoria
acumulada es dinámica y trazable.

\section{Calendario dodecafásico y no escisión}

El calendario interno se organiza mediante
\[
 0\longrightarrow C_{12}\longrightarrow C_{108}
 \longrightarrow C_9\longrightarrow0.
\]
La extensión conserva la diferencia entre una vuelta visible y el estado
levantado. En la realización
\(C_{108}\simeq C_{27}\times C_4\), la parte \(C_{27}\to C_9\) no se
escinde: los levantamientos del generador tienen orden \(27\), no orden
divisor de \(9\). El factor \(C_4\) aporta el cuarto de giro; el núcleo
ternario aporta la elevación. La combinación no es una etiqueta temporal
externa, sino la arquitectura de retorno del TPK.

\section{Posterioridad del cociente de Markov como publicación de memoria reducida}

Promediar una dinámica enriquecida puede producir un operador de memoria
reducida \(K_{\ph}\). Este operador es útil para estudiar transferencia,
estacionariedad y espectro, pero no genera \(\Gamma_9\). La dependencia es
\[
 U_t\longrightarrow\Gamma_9\longrightarrow K_{\ph}.
\]
Invertirla borraría hoja, prefijo, frontera y carry. En esta publicación el
lector radial se construye sobre las historias enriquecidas, nunca sobre el
cociente de Markov.

\part{Geometría helicoidal exacta de la memoria nonádica}

\chapter{Hélice nonádica de fase y profundidad}

\section{Levantamiento de la fase nonádica al recubrimiento fase--memoria}

Sea
\[
 \operatorname{Rep}_9=\{0,1,\ldots,8\},
 \qquad \widetilde\Theta=\operatorname{Rep}_9\times\Z.
\]
Definimos
\begin{equation}
 \widetilde\Gamma_9(r,n)=(r,n+1),
 \qquad
 \Pi(r,n)=[r+9n]_{108}.
 \label{espiral:eq:cobertura-helice}
\end{equation}
La primera coordenada registra la fase visible y la segunda la profundidad.
Se cumplen
\[
 \Pi(r,n+12)=\Pi(r,n),
 \qquad
 \widetilde\Gamma_9^{12}(r,n)=(r,n+12).
\]
La transformación de cubierta
\[
 \tau_{12}(r,n)=(r,n+12)
\]
actúa libremente. Cada fibra de \(\Pi\) es una órbita de esa acción.

\begin{definition}[Hélice nonádica]
La hélice nonádica es el sistema levantado
\[
 (\widetilde\Theta,\widetilde\Gamma_9,\Pi)
\]
con fase \(r\bmod9\), calendario \(\Pi(r,n)\) y memoria profunda \(n\).
\end{definition}

\begin{theorem}[Fase periódica y memoria no periódica]
La proyección de la hélice es periódica de orden \(12\) sobre el calendario
\(C_{108}\), mientras el estado levantado distingue todas las profundidades
\(n+12k\), \(k\in\Z\).
\end{theorem}
\begin{proof}
La igualdad \(\Pi(r,n)=\Pi(r',n')\) obliga a \(r=r'\) y
\(n-n'\in12\Z\). Por tanto la proyección identifica exactamente las
órbitas de \(\tau_{12}\), pero la segunda coordenada del levantamiento
permanece entera y distingue sus elementos.
\end{proof}

\section{Inmersión geométrica normalizada}

Para \(N\ge1\), sea
\[
 z_N=q_9(N)+\frac{\rho_9(N)-1}{9}.
\]
La inmersión
\begin{equation}
 H_9(N)=\bigl(\cos(2\pi z_N),\sin(2\pi z_N),z_N\bigr)
 \label{espiral:eq:inmersion-helice-nonadica}
\end{equation}
satisface
\[
 z_{N+9}=z_N+1,
 \qquad
 H_9(N+9)=H_9(N)+(0,0,1).
\]
La proyección sobre las dos primeras coordenadas retorna; la tercera conserva
la vuelta. La figura no se impone a la aritmética: es una inmersión de la
descomposición exacta \(N=9q_9(N)+\rho_9(N)\).

Para la curva continua
\[
 h(z)=(\cos2\pi z,\sin2\pi z,z),
\]
se tiene
\[
 h'(z)=(-2\pi\sin2\pi z,2\pi\cos2\pi z,1),
\]
\[
 h''(z)=(-4\pi^2\cos2\pi z,-4\pi^2\sin2\pi z,0).
\]
Su curvatura y torsión son constantes:
\begin{equation}
 \kappa_h=\frac{4\pi^2}{4\pi^2+1},
 \qquad
 \mathcal T_h=\frac{2\pi}{4\pi^2+1}.
 \label{espiral:eq:frenet-helice-normalizada}
\end{equation}

\begin{proof}[Cálculo de \cref{espiral:eq:frenet-helice-normalizada}]
La norma de \(h'\) es \(\sqrt{4\pi^2+1}\) y
\(\|h'\times h''\|=4\pi^2\sqrt{4\pi^2+1}\). Por tanto
\[
 \kappa_h=\frac{\|h'\times h''\|}{\|h'\|^3}
 =\frac{4\pi^2}{4\pi^2+1}.
\]
Además, con \(h'''=(8\pi^3\sin2\pi z,-8\pi^3\cos2\pi z,0)\),
\[
 \mathcal T_h=
 \frac{\det(h',h'',h''')}{\|h'\times h''\|^2}
 =\frac{2\pi}{4\pi^2+1}.
\]
\end{proof}

\begin{figure}[htbp]
\centering
\begin{tikzpicture}[x={(1cm,0cm)},y={(.30cm,.12cm)},z={(0cm,.68cm)},scale=.82]
 \draw[->,HMTGray] (0,0,0)--(0,0,7.2) node[above] {profundidad};
 \foreach \k in {0,...,4}{
   \draw[HMTLine] (0,0,{1.4*\k}) ellipse (1.55 and .48);
 }
 \draw[HMTBlue,very thick,domain=0:1440,samples=260,smooth,variable=\t]
   plot ({1.55*cos(\t)},{1.55*sin(\t)},{\t/210});
 \foreach \j in {0,...,35}{
   \fill[HMTRed] ({1.55*cos(40*\j)},{1.55*sin(40*\j)},{40*\j/210}) circle (1.1pt);
 }
 \node[HMTGray,font=\small\sffamily,right] at (1.7,0,3.3) {$9$ marcas por vuelta};
\end{tikzpicture}
\caption{Hélice nonádica: los anillos son secciones de nivel; la curva conserva el avance de memoria.}
\label{espiral:fig:helice-nonadica-analitica}
\end{figure}

\section{Círculos anidados y cilindros}

Una sección \(n=\text{constante}\) produce un círculo. Una familia de
secciones con radio constante produce un cilindro. Si el radio depende de la
profundidad, la misma estructura produce una suspensión helicoidal radial.
Por tanto:
\[
 \text{órbita}\longleftrightarrow\text{hoja helicoidal},
\]
\[
 \text{sección de nivel}\longleftrightarrow\text{círculo},
\qquad
 \text{familia de secciones}\longleftrightarrow\text{cilindro}.
\]
Ésta es una traducción geométrica del sistema levantado, no una sustitución
del ledger por un dibujo.

\chapter{Suspensión dinámica del sistema afín enriquecido}

\section{Holonomía dependiente del estado}

Sea \(U_{\APP}\) un módulo radial y asóciese a cada bloque nonádico un
endomorfismo afín
\[
 H_9(x):u\longmapsto\Lambda_9(x)u+C_9(x).
\]
Cuando la holonomía depende del estado de partida, definimos
\(F(x,u)=(\Gamma_9x,H_9(x)u)\). Si \(F\) es continua, su suspensión es
el cociente de mapeo
\begin{equation}
 \mathfrak S_{\TPK}^{\rad}
 =\frac{X_{\TPK}^{\enr}\times U_{\APP}\times[0,1]}
 {(x,u,1)\sim(\Gamma_9x,H_9(x)u,0)}.
 \label{espiral:eq:suspension-skew}
\end{equation}
Si \(F\) es un homeomorfismo ---lo cual exige, además de la invertibilidad
de las partes lineales, la reversibilidad continua de \(\Gamma_9\) y de su
dependencia en \(x\)--- se obtiene una suspensión bidireccional. En el
monoide general el mismo cociente codifica sólo la evolución hacia delante;
no se lo identifica con un telescopio ni se presupone una inversión global.

\begin{definition}[Suspensión helicoidal radial]
La suspensión \cref{espiral:eq:suspension-skew}, provista de una coordenada angular
de fase y una evaluación radial posterior, se denomina \emph{suspensión
helicoidal radial}. El término \emph{hélice general} se reserva para los
casos que satisfacen una condición diferencial de Lancret.
\end{definition}

\section{Hojas orbitales y costura del espacio de suspensión}

La órbita del skew--product es
\[
 (x_0,u_0),
 (\Gamma_9x_0,H_9(x_0)u_0),
 (\Gamma_9^2x_0,H_9(\Gamma_9x_0)H_9(x_0)u_0),\ldots.
\]
La costura conserva el orden de los factores. Sustituir toda la órbita por el
producto final impediría truncar la historia y reconstruir sus prefijos. La
hoja helicoidal es, por tanto, una órbita con ledger, no una curva aislada.

\section{Caso estacionario afín}

Si \(H_9(x)=H=(\Lambda,C)\) es constante,
\[
 u_{n+1}=\Lambda u_n+C.
\]
La iteración interna exacta, válida para todo endomorfismo \(\Lambda\), es
\begin{equation}
 u_n=\Lambda^n u_0+\sum_{k=0}^{n-1}\Lambda^kC.
 \label{espiral:eq:orbita-afin-estacionaria-general}
\end{equation}
Si \(I-\Lambda\) es invertible en el módulo radial, existe el punto fijo
\[
 u_*=(I-\Lambda)^{-1}C,
\]
y
\begin{equation}
 u_n-u_*=\Lambda^n(u_0-u_*).
 \label{espiral:eq:orbita-afin-estacionaria}
\end{equation}
En el sector traslacional \(\Lambda=I\),
\[
 u_n=u_0+nC.
\]
Sólo después de elegir un carácter escalar positivo
\(\chi:U_{\APP}\to\R\), con
\(\chi\Lambda=\lambda\chi\), se obtienen las fórmulas numéricas
\(u_*=C/(1-\lambda)\) para \(\lambda\neq1\) y las potencias reales de
\(\lambda\) empleadas en la interpolación continua. No se divide un
endomorfismo abstracto por \(1-\Lambda\).
Estas dos posibilidades anticipan la separación entre una familia de
autoescala y una familia traslacional. Ninguna selecciona por sí sola una
espiral conocida; la curva aparece después de aplicar \(\exp_p\) y la
coordenada angular.

\begin{figure}[htbp]
\centering
\begin{tikzpicture}[>=Latex,node distance=9mm and 14mm,
 s/.style={draw=HMTNavy,fill=white,minimum width=21mm,minimum height=9mm,font=\small\sffamily}]
\node[s] (x0) {\((x_0,u_0)\)};
\node[s,right=of x0] (x1) {\((x_1,u_1)\)};
\node[s,right=of x1] (x2) {\((x_2,u_2)\)};
\node[right=7mm of x2] (dots) {$\cdots$};
\draw[->,HMTBlue,thick] (x0)--node[above,font=\scriptsize] {$\Gamma_9,H_9(x_0)$} (x1);
\draw[->,HMTBlue,thick] (x1)--node[above,font=\scriptsize] {$\Gamma_9,H_9(x_1)$} (x2);
\draw[->,HMTBlue,thick] (x2)--(dots);
\draw[HMTRed,dashed,rounded corners] ([xshift=-3mm,yshift=-4mm]x0.south west) rectangle ([xshift=3mm,yshift=4mm]x2.north east);
\node[below=7mm of x1,HMTGray,font=\small\sffamily] {cada prefijo permanece accesible};
\end{tikzpicture}
\caption{La suspensión conserva la secuencia de factores, no sólo la holonomía final.}
\label{espiral:fig:skew-ledger}
\end{figure}

\chapter{Involución bidireccional y hojas conjugadas}

\section{Inversión de rutas}

Sea \(\mathscr P_{\TPK}^{\times}\) el subgrupoide de rutas reversibles. Si
la representación radial satisface
\[
 h_{e^\dagger}=h_e^{-1},
\]
entonces, para
\(\gamma=e_n\cdots e_1\),
\[
 H_{C_{\mathrm{rev}}\gamma}
 =h_{e_1}^{-1}\cdots h_{e_n}^{-1}
 =(h_{e_n}\cdots h_{e_1})^{-1}
 =H_\gamma^{-1}.
\]

\begin{theorem}[Conjugación de ascenso y descenso]
En el sector reversible, la involución de rutas intercambia las hojas de
ascenso y descenso e invierte la holonomía afín.
\end{theorem}

Sobre la cobertura helicoidal, sea
\[
 \iota(r,n)=(r,-n).
\]
Entonces
\begin{equation}
 \iota\widetilde\Gamma_9\iota
 =\widetilde\Gamma_9^{-1}.
 \label{espiral:eq:involucion-helice}
\end{equation}
La igualdad se obtiene directamente:
\[
 (r,n)\xmapsto{\iota}(r,-n)
 \xmapsto{\widetilde\Gamma_9}(r,-n+1)
 \xmapsto{\iota}(r,n-1).
\]

\begin{figure}[htbp]
\centering
\begin{tikzpicture}[x={(1cm,0cm)},y={(.30cm,.12cm)},z={(0cm,.58cm)},scale=.72]
 \draw[->,HMTGray] (0,0,-4.8)--(0,0,4.8) node[above] {memoria};
 \draw[HMTBlue,very thick,domain=0:900,samples=180,smooth,variable=\t]
   plot ({1.25*cos(\t)},{1.25*sin(\t)},{\t/190});
 \draw[HMTRed,very thick,domain=0:900,samples=180,smooth,variable=\t]
   plot ({1.25*cos(-\t)},{1.25*sin(-\t)},{-\t/190});
 \fill[HMTNavy] (1.25,0,0) circle (1.5pt);
 \node[right,HMTBlue,font=\small\sffamily] at (1.5,0,3.4) {ascenso};
 \node[right,HMTRed,font=\small\sffamily] at (1.5,0,-3.4) {descenso};
\end{tikzpicture}
\caption{Hojas conjugadas producidas por la inversión de ruta.}
\label{espiral:fig:hojas-conjugadas}
\end{figure}

\section{Espejo de constantes y eje fijo}

La involución sobre la fibra de constantes
\[
 \mathfrak c(u^\pi,u^e,u^\varphi)
 =(u^e,u^\pi,u^\varphi)
\]
intercambia las ramas \(\pi/e\) y fija el eje \(\varphi\). También conserva
\(u^\pi+u^e\). Esta involución y la inversión de rutas son objetos distintos:
la primera actúa sobre una fibra de publicación; la segunda actúa sobre el
grupoide de caminos. No se afirma una conjugación global
\(\mathfrak c\Gamma_9\mathfrak c=\Gamma_9^{-1}\) sin construir el
entrelazador correspondiente.

Esta distinción protege la tesis positiva. El espejo \(\pi/e\), el eje
\(\varphi\) y la bidireccionalidad están presentes; lo que no se permite es
identificarlos por una igualdad de operadores no tipada.

\section{Círculos anidados como cortes de las hojas}

Sea \(r_n>0\) la escala de la sección \(n\). La familia
\[
 C_n=\{(r_n\cos\theta,r_n\sin\theta,n):0\le\theta<2\pi\}
\]
produce círculos anidados cuando \(r_n\) varía y cilindros cuando se agrupan
secciones con radio común. La unión de los \(C_n\) a lo largo de una ruta
TPK es la traza discreta de una hoja helicoidal. La involución
\(n\mapsto-n\) produce la hoja conjugada.
